\begin{theorem}[Triangle Inequality]
Let $(V,\langle\cdot,\cdot\rangle)$ be an inner‑product space and let $\|\cdot\|$ be the norm induced by the inner product, $\|z\|=\sqrt{\langle z,z\rangle}$.  
For any $x,y\in V$,
\[
\|x+y\|\le \|x\|+\|y\|.
\]
\end{theorem}

\begin{proof}
\textbf{1. Degenerate cases.}  
If $x=0$ then $\|x+y\|=\|y\|$ and the right‑hand side is $\|0\|+\|y\|=\|y\|$; similarly for $y=0$. Hence the inequality holds with equality. In the sequel we assume $x\neq0$ and $y\neq0$.

\medskip\noindent\textbf{2. Squaring.}  
Both $\|x+y\|$ and $\|x\|+\|y\|$ are non‑negative, therefore
\begin{equation}
\|x+y\|\le\|x\|+\|y\|\iff
\|x+y\|^{2}\le(\|x\|+\|y\|)^{2}.
\tag{1}
\end{equation}

\medskip\noindent\textbf{3. Expansion of the left‑hand side.}  
Using linearity in the first argument and conjugate symmetry,
\begin{align*}
\|x+y\|^{2}
&=\langle x+y,\,x+y\rangle\\
&=\langle x,x\rangle+\langle x,y\rangle+\langle y,x\rangle+\langle y,y\rangle\\
&=\|x\|^{2}+\|y\|^{2}+ \langle x,y\rangle+\overline{\langle x,y\rangle}\\
&=\|x\|^{2}+\|y\|^{2}+2\operatorname{Re}\langle x,y\rangle .
\end{align*}
\tag{2}
\]

The right‑hand side of \eqref{1} expands as
\begin{equation}
(\|x\|+\|y\|)^{2}= \|x\|^{2}+2\|x\|\|y\|+\|y\|^{2}.
\tag{3}
\end{equation}

\medskip\noindent\textbf{4. Bounding the mixed term.}  
By the Cauchy–Schwarz inequality,
\begin{equation}
|\langle x,y\rangle|\le\|x\|\,\|y\|.
\tag{4}
\end{equation}
Since $\operatorname{Re}z\le|z|$ for any $z\in\mathbb{C}$, we obtain
\begin{equation}
\operatorname{Re}\langle x,y\rangle\le|\langle x,y\rangle|\le\|x\|\,\|y\|.
\tag{5}
\end{equation}

\medskip\noindent\textbf{5. Comparison.}  
Substituting \eqref{5} into \eqref{2} gives
\[
\|x+y\|^{2}
=\|x\|^{2}+2\operatorname{Re}\langle x,y\rangle+\|y\|^{2}
\le\|x\|^{2}+2\|x\|\|y\|+\|y\|^{2}
=(\|x\|+\|y\|)^{2}.
\]
Thus \eqref{1} holds, and taking square roots yields $\|x+y\|\le\|x\|+\|y\|$ for all $x,y\in V$.

\medskip\noindent\textbf{6. Equality condition.}  
Equality occurs precisely when equality holds in both \eqref{4} and in $\operatorname{Re}\langle x,y\rangle\le|\langle x,y\rangle|$.

Equality in \eqref{4} (Cauchy–Schwarz) means $x$ and $y$ are linearly dependent, i.e. $x=\lambda y$ for some $\lambda\in\mathbb{F}$ (or $y=0$).  
Equality in $\operatorname{Re}\langle x,y\rangle\le|\langle x,y\rangle|$ forces $\langle x,y\rangle$ to be a non‑negative real number.  
Hence $\|x+y\|=\|x\|+\|y\|$ iff either $x=0$, $y=0$, or $x=\alpha y$ with $\alpha\ge0$ (real). In other words, the two vectors point in the same direction (or one is the zero vector).
\end{proof}